\chapter{\texorpdfstring{$t$}{t}分布}
    \section{定義}
        \begin{shadebox}
            関数$f(t)$を次のように定義する。
            \[f(t) \coloneqq \frac{\Gamma{\left(\frac{n+1}{2}\right)}}{\sqrt{n\pi}\Gamma{\left(\frac{n}{2}\right)}}\left(1+\frac{t^2}{n}\right)^{-\frac{n+1}{2}}\]
            確率密度関数が上式で与えられるような分布を「自由度$n$の$t$分布」と呼び、$t_n$と書くことにする。
        \end{shadebox}
    \section{\thesubsection \texorpdfstring{$X,Y : \text{indep},\;\;\;X \sim N(0,1),\;\;\;Y \sim {\chi^2}_n \quad \Rightarrow \quad T \coloneqq \frac{X}{\sqrt{\frac{Y}{2}}} \sim t_n$}{X,Y: 独立, X～N(0,1), Y～χ\string^2\_n ⇒ T := X/√(Y/2) ～ t\_n}}
        \begin{shadebox}
            \[X,Y : \text{indep},\;\;\;X \sim N(0,1),\;\;\;Y \sim {\chi^2}_n \quad \Rightarrow \quad T \coloneqq \frac{X}{\sqrt{\frac{Y}{2}}} \sim t_n\]
        \end{shadebox}
        \begin{proof}
            \quad\par
            $X,Y,T$の確率密度関数を$f^X(x),f^Y(y),f~T(t)$と書くことにする。
            \begin{align*}
                &\quad P(-\infty \leq T \leq t) = P\left(-\infty \leq X \leq t\sqrt{\frac{y}{n}}\right) = \int_{0}^{\infty}{\int_{-\infty}^{t\sqrt{\frac{y}{n}}}{f^Y(y)f^X(x) dx} dy} \\
                &= \int_{0}^{\infty}{f^Y(y)\left(\int_{-\infty}^{t\sqrt{\frac{y}{n}}}{f^X(x) dx}\right) dy}
            \end{align*}
            \begin{align*}
                \therefore \; f^T(t) &= \frac{\partial}{\partial t}P(-\infty \leq T \leq t) = \int_{0}^{\infty}{f^Y(y)f^X\left(t\sqrt{\frac{y}{n}}\right)\sqrt{\frac{y}{n}} dy}\quad\text{(微分と積分の順序を交換した)}\\
                &=\frac{1}{\sqrt{2n\pi}2^{\frac{n}{2}}\Gamma{\left(\frac{n}{2}\right)}}\int_{0}^{\infty}{y^{\frac{n}{2}-1}e^{-\frac{y}{2}}\exp{\left(-\frac{t^2y}{2n}\right)}y^{\frac{1}{2}} dy} \\
                &= \frac{1}{\sqrt{2n\pi}2^{\frac{n}{2}}\Gamma{\left(\frac{n}{2}\right)}}\int_{0}^{\infty}{y^{\frac{n-1}{2}}\exp{\left[-\frac{y}{2}\left(1+\frac{t^2}{n}\right)\right]} dy}\\
                &\left(\frac{y}{2}\left(1+\frac{t^2}{n}\right) = u \text{なる変数変換を行って}\right)\\
                &=\frac{1}{\sqrt{2n\pi}2^{\frac{n}{2}}\Gamma{\left(\frac{n}{2}\right)}}\left(\frac{2}{1+\frac{t^2}{n}}\right)^{\frac{n-1}{2}}\int_{0}^{\infty}{u^{\frac{n-1}{2}}e^{-u}\frac{2}{1+\frac{t^2}{n}} du}\\
                &=\frac{1}{\sqrt{2n\pi}2^{\frac{n}{2}}\Gamma{\left(\frac{n}{2}\right)}}\left(\frac{2}{1+\frac{t^2}{n}}\right)^{\frac{n+1}{2}}\int_{0}^{\infty}{u^{\frac{n+1}{2}-1}e^{-u} du} = \frac{1}{\sqrt{2n\pi}2^{\frac{n}{2}}\Gamma{\left(\frac{n}{2}\right)}}\left(\frac{2}{1+\frac{t^2}{n}}\right)^{\frac{n+1}{2}}\Gamma{\left(\frac{n+1}{2}\right)}\\
                &= \frac{\Gamma{\left(\frac{n+1}{2}\right)}}{\sqrt{n\pi}\Gamma{\left(\frac{n}{2}\right)}}\left(1+\frac{t^2}{n}\right)^{-\frac{n+1}{2}} = \text{pdf of} \quad t_n
            \end{align*}
        \end{proof}
    \section{\thesubsection \texorpdfstring{$t$}{t}分布の期待値}
        \begin{shadebox}
            $t$分布は自由度$n \geq 2$の場合に限り期待値が存在し$E[T] = 0$である。
        \end{shadebox}
        \begin{proof}
            \quad\par
            $tf(t)$が$t \to \infty$で$\frac{1}{t}$オーダーより小さくなる条件が$n \geq 2$であることは式を見れば容易にわかる。
            この時に限り$\int_{0}^{\infty}{tf(t) dt}$および$\int_{-\infty}^{0}{tf(t) dt}$が存在し$E[T] = \int_{-\infty}^{\infty}{tf(t) dt}$が存在する。
            被積分関数が奇関数であることから積分値は明らかに0である。
        \end{proof}
        \section{\thesubsection \texorpdfstring{$t$}{t}分布の分散}
        \begin{shadebox}
            $t$分布の分散は自由度$n \leq 2$の場合$\infty$、自由度$n \geq 3$の場合$\frac{n}{n-2}$
        \end{shadebox}
        \begin{proof}
            \quad\par
            \setcounter{equation}{0}
            $n \leq 2$の場合、$t \to\pm\infty$で$t^2f(t)$が$\frac{1}{t}$オーダーを下回らないから$\int_{-\infty}^{\infty}{t^2f(t) dt} = \infty$ \\
            $n \geq 3$の場合、積分は収束する。このとき分散が$\frac{n}{n-2}$であることを示す。\\
            確率密度関数は偶関数だから
            \begin{equation}
                \int_{-\infty}^{\infty}{t^2f(t) dt} = 2\int_{0}^{\infty}{t^2f(t) dt} = \frac{2\Gamma{\left(\frac{n+1}{2}\right)}}{\sqrt{n\pi}\Gamma{\left(\frac{n}{2}\right)}}\int_{0}^{\infty}{t^2\left(1+\frac{t^2}{n}\right)^{-\frac{n+1}{2}} dt}
            \end{equation}
            先に積分の値を計算する。$t^2 = u$なる変数変換を行って
            \[\int_{0}^{\infty}{t^2\left(1+\frac{t^2}{n}\right)^{-\frac{n+1}{2}} dt} = \int_{0}^{\infty}{u\left(1+\frac{u}{n}\right)^{-\frac{n+1}{2}} \frac{1}{2\sqrt{u}}du} = \frac{1}{2}\int_{0}^{\infty}{u^{\frac{1}{2}}\left(1+\frac{u}{n}\right)^{-\frac{n+1}{2}} du}\]
            次に$\displaystyle v = \frac{\frac{u}{n}}{1+\frac{u}{n}} = \frac{u}{n+u}$なる変数変換を行う。$\displaystyle u = \frac{nv}{1-v}$となって$du = \frac{n}{(1-v)^2}dv$となる。積分範囲は$v=0 \to \textcolor{red}{1}$になる。\\
            (※(1)で積分の左側に括りだした定数群にガンマ関数が居る。最終的にこれらが消えてほしいわけだから積分結果がガンマ関数あるいはベータ関数にならなければならない。そこでベータ関数の出現を期待して積分範囲を$0\to\infty$から$0\to 1$に変更するような変数変換を行ったのである。)\\
            上式は
            \begin{align*}
                & \frac{1}{2}\int_{0}^{1}{\left(\frac{nv}{1-v}\right)^{\frac{1}{2}}\left(1+\frac{v}{1-v}\right)^{-\frac{n+1}{2}}\frac{n}{(1-v)^2} dv} = \frac{n\sqrt{n}}{2}\int_{0}^{1}{\sqrt{v}(1-v)^{\frac{n}{2}-2} dv}\\
                &= \frac{n\sqrt{n}}{2}\int_{0}^{1}{v^{\frac{3}{2}-1}(1-v)^{\frac{n-2}{2}-1} dv} = \frac{n\sqrt{n}}{2}B\left(\frac{3}{2},\frac{n-2}{2}\right) = \frac{n\sqrt{n}}{2}\frac{\Gamma{\left(\frac{3}{2}\right)}\Gamma{\left(\frac{n-2}{2}\right)}}{\Gamma{\left(\frac{n+1}{2}\right)}}
            \end{align*}
            ここで
            \[\Gamma{\left(\frac{3}{2}\right)} = \Gamma{\left(\frac{1}{2}+1\right)} = \frac{1}{2}\Gamma{\left(\frac{1}{2}\right)} = \frac{\sqrt{\pi}}{2}\]
            および
            \[\Gamma{\left(\frac{n-2}{2}\right)} = \Gamma{\left(\frac{n}{2}-1\right)}\]
            であり
            \[\Gamma{\left(\frac{n}{2}\right)} = \Gamma{\left(\frac{n}{2}-1+1\right)} = \left(\frac{n}{2}-1\right)\Gamma{\left(\frac{n}{2}-1\right)}\]
            であるから
            \[\Gamma{\left(\frac{n}{2}-1\right)} = \frac{2\Gamma{\left(\frac{n}{2}\right)}}{n-2}\]
            以上の結果を(1)に適用して
            \[\int_{-\infty}^{\infty}{t^2f(t) dt} = \frac{n}{n-2}\]
            最後に
            \[V[T] = E[T^2] - {E[T]}^2 = \int_{-\infty}^{\infty}{t^2f(t) dt} - 0 = \frac{n}{n-2}\]
        \end{proof}